% ---------------------------------------------------------------------
\section{Les règles du jeu en France}
% ---------------------------------------------------------------------

Un livre imprimé à la demande est un livre comme un autre : toutes les
règles du livre en France s’appliquent, que vous soyez une maison
d’édition ou un auteur qui s’autoédite.

\subsection{Le prix unique du livre (loi Lang, 1981)}

\begin{itemize}
  \item \textbf{C’est l’éditeur qui fixe le prix public} TTC, et ce prix
        s’impose à tous les vendeurs en France : Amazon, Fnac, libraire
        de quartier, salon du livre.
  \item Le vendeur peut accorder \textbf{au maximum 5 \% de remise}
        sur ce prix.
  \item \textbf{Le prix doit être imprimé} sur le livre (en général en
        quatrième de couverture).
  \item Si vous vous autoéditez, \textbf{vous êtes l’éditeur} : vous
        fixez donc le prix… et vous devez le respecter vous-même, y
        compris en vente directe.
\end{itemize}

\begin{attention}[Plateformes étrangères]
Certaines plateformes calculent un prix « local » à partir d’un prix en
dollars ou en euros saisi ailleurs. Vérifiez que le prix affiché en
France est \textbf{identique sur tous les canaux} et correspond à celui
imprimé sur le livre.
\end{attention}

\subsection{Les frais de port (lois de 2014 et loi Darcos de 2021)}

Pour protéger les librairies, la vente de livres neufs en ligne avec
livraison à domicile est encadrée :

\chiffrecle{3~€}{Tarif minimal TTC de livraison d’un livre neuf pour une
commande de moins de 35~€ (arrêté d’application de la loi Darcos, en
vigueur depuis octobre 2023). Au-delà de 35~€, le minimum est symbolique
(0,01~€). Le retrait en magasin peut rester gratuit.}

Cela concerne aussi \textbf{votre propre boutique en ligne} si vous
vendez vos livres neufs vous-même.

\subsection{La TVA}

Le livre papier bénéficie du \textbf{taux réduit de 5,5 \%} en France
métropolitaine. Les redevances des plateformes sont calculées sur le
prix \textbf{hors taxe}. Un livre vendu \eur{\exPrixTTC} TTC rapporte donc
\eur{\exPrixHT} HT avant toute autre déduction.

\subsection{L’ISBN}

\begin{itemize}
  \item L’ISBN identifie une \textbf{édition} d’un livre :
        il en faut un par format (broché, relié, numérique, gros
        caractères…). Changer le prix ne nécessite pas de nouvel ISBN ;
        changer le format ou le contenu de façon importante, si.
  \item En France, il est attribué \textbf{gratuitement par l’AFNIL}
        (\href{https://www.afnil.org}{afnil.org}), à l’éditeur.
        Un auteur qui s’autoédite peut en demander à son nom.
  \item Beaucoup de plateformes offrent un ISBN gratuit, mais
        \textbf{la plateforme apparaît alors comme éditeur} (ou le livre
        est enregistré sous une mention du type « Independently
        published »). Vous ne pourrez pas transférer ce livre et cet
        ISBN ailleurs.
\end{itemize}

\begin{astuce}[Notre conseil]
Prenez \textbf{vos propres ISBN} auprès de l’AFNIL : c’est gratuit,
vous restez l’éditeur officiel, et vous pouvez changer de plateforme ou
en utiliser plusieurs sans repartir de zéro.
\end{astuce}

\subsection{Le dépôt légal}

Le dépôt légal à la \textbf{Bibliothèque nationale de France} est
\textbf{obligatoire} pour tout livre diffusé en France, y compris en IAD
et en autoédition. C’est à l’\textbf{éditeur} (donc à vous si vous vous
autoéditez) de déposer le ou les exemplaires demandés, accompagnés d’une
déclaration, au plus tard le jour de la mise en vente. L’envoi à la BnF
est gratuit (franchise postale).

\source{Le nombre d’exemplaires et la procédure (déclaration en ligne
notamment) sont détaillés sur bnf.fr, rubrique « Dépôt légal ». Les
réimpressions à l’identique d’un livre en IAD n’ont pas à être déposées
de nouveau ; une nouvelle édition, oui.}

\subsection{Les mentions obligatoires}

Sur la page de copyright (en général au verso de la page de titre) ou
en fin d’ouvrage, on trouve :

\begin{itemize}
  \item le \textbf{nom et l’adresse de l’éditeur} ;
  \item le \textbf{nom et l’adresse de l’imprimeur} (en IAD : la mention
        fournie par la plateforme, par exemple « Imprimé à la demande
        par… ») ;
  \item la \textbf{date d’achèvement du tirage} (en IAD, on indique
        souvent la date de la première mise en impression) ;
  \item la \textbf{mention « Dépôt légal : mois année »} ;
  \item l’\textbf{ISBN} et le \textbf{prix} ;
  \item pour les livres destinés à la jeunesse : la mention de la
        \textbf{loi n° 49-956 du 16 juillet 1949} sur les publications
        destinées à la jeunesse.
\end{itemize}

\subsection{Sécurité des produits (règlement européen GPSR)}

Depuis le 13 décembre 2024, le \textbf{règlement (UE) 2023/988 sur la
sécurité générale des produits} s’applique aussi aux livres vendus dans
l’Union. En pratique, les plateformes demandent désormais le nom et
l’adresse (postale et électronique) d’un \textbf{responsable établi dans
l’UE} — l’éditeur, en général — pour chaque titre. Prévoyez ces
informations au moment de la mise en ligne.

\subsection{Statut : auteur, auteur-éditeur ou maison d’édition}

\subsubsection{Vous êtes auteur et vous vous autoéditez}

Les revenus tirés de la vente de vos propres livres en autoédition
peuvent, sous conditions, relever du \textbf{régime des
artistes-auteurs} (déclaration à l’Urssaf Limousin, service
artistes-auteurs). Selon le volume et la nature de l’activité (achat et
revente de stocks, vente sur salons…), une \textbf{micro-entreprise}
peut être plus adaptée. Faites le point avec l’Urssaf ou un comptable :
c’est le sujet où les erreurs coûtent le plus cher.

\subsubsection{Vous créez ou dirigez une petite maison d’édition}

\begin{itemize}
  \item Forme juridique : association, micro-entreprise, SAS ou SARL
        (code APE habituel : \textbf{58.11Z — Édition de livres}).
  \item Avec vos auteurs, vous signez un \textbf{contrat d’édition}
        (Code de la propriété intellectuelle, art. L132-1 et suivants),
        avec une \textbf{reddition des comptes} au moins annuelle et une
        partie distincte pour l’édition numérique.
  \item L’éditeur doit assurer une \textbf{exploitation permanente et
        suivie} de l’œuvre. L’IAD aide à garder le livre disponible,
        mais il doit rester présenté au catalogue et réellement
        commandable ; vérifiez ce que prévoient vos contrats.
\end{itemize}

\subsection{À surveiller}

\begin{itemize}
  \item \textbf{Le règlement européen contre la déforestation (EUDR)},
        qui concerne le papier imprimé : son entrée en application a été
        repoussée à plusieurs reprises. Les plateformes et imprimeurs
        devraient gérer l’essentiel, mais on pourra vous demander des
        informations.
  \item \textbf{Les débats sur le livre d’occasion} et la rémunération
        des auteurs, régulièrement relancés au Parlement et au ministère
        de la Culture.
\end{itemize}
